\documentclass[a4paper,12pt]{article}

\usepackage{iftex}
\ifPDFTeX
  \usepackage[T2A]{fontenc}
  \usepackage[utf8]{inputenc}
  \usepackage[russian]{babel}
\else
  \usepackage{fontspec}
  \usepackage{polyglossia}
  \setmainlanguage{russian}
  \setmainfont{DejaVu Sans}
\fi
\usepackage{amsmath,amssymb,amsthm,mathtools}
\ifPDFTeX
  \usepackage{helvet}
  \renewcommand{\familydefault}{\sfdefault}
\fi
\usepackage{icomma}
\usepackage[left=18mm,right=18mm,top=18mm,bottom=20mm]{geometry}
\usepackage{microtype}
\usepackage{tikz}
\usepackage{tkz-euclide}
\usepackage{pgfplots}
\pgfplotsset{compat=1.18}
\usepackage{tabularx,booktabs,longtable}
\usepackage{enumitem}
\usepackage{hyperref}
\usepackage{parskip}
\usepackage{fancyhdr}
\usepackage{setspace}
\usepackage{xcolor}

\definecolor{accent}{HTML}{008080}
\definecolor{darkaccent}{HTML}{004D4D}
\definecolor{textgray}{HTML}{333333}
\definecolor{softgray}{HTML}{666666}

\onehalfspacing
\hypersetup{colorlinks=true,linkcolor=darkaccent,urlcolor=darkaccent}
\pagestyle{fancy}
\fancyhf{}
\fancyhead[L]{\small Чек-лист: прикладные формулы}
\fancyhead[R]{\small Ксения Клыкова}
\fancyfoot[C]{\small \thepage}
\renewcommand{\headrulewidth}{0.4pt}
\setlength{\parindent}{0pt}
\setlength{\parskip}{6pt}
\setlength{\headheight}{14pt}
\setlength{\emergencystretch}{3em}
\sloppy
\setlist[itemize]{leftmargin=1.15cm,itemsep=2pt,topsep=3pt}
\setlist[enumerate]{leftmargin=1.15cm,itemsep=3pt,topsep=3pt}
\newcommand{\sectiontitle}[1]{\section*{\textcolor{darkaccent}{#1}}\addcontentsline{toc}{section}{#1}}
\newcommand{\subsectiontitle}[1]{\subsection*{\textcolor{accent}{#1}}}
\newcommand{\important}[1]{\textbf{\textcolor{darkaccent}{#1}}}

\begin{document}

\begin{titlepage}
\thispagestyle{empty}
\begin{center}
\vspace*{27mm}
{\Large \textcolor{softgray}{Чек-лист}}
\vspace{9mm}

{\Huge \textbf{Прикладные формулы}}

\vspace{6mm}
{\Large \textbf{Перевод условия, зависимости, степени и проверка ответа}}

\vspace{22mm}
{\large Лёвин Артём Александрович эксклюзивно для Ксении Клыковой}

\vspace{12mm}
{\large 15 июля 2026 года}

\vfill
\begin{tikzpicture}
\draw[accent,line width=1.1pt]
  (0,0) .. controls (2,0.85) and (4,-0.85) .. (6,0)
        .. controls (8,0.85) and (10,-0.85) .. (12,0);
\end{tikzpicture}
\vspace{14mm}
\end{center}
\end{titlepage}

\tableofcontents
\newpage

\sectiontitle{1. Универсальная стратегия}

Прикладная задача проверяет умение перевести реальную ситуацию на язык математики и аккуратно преобразовать формулу. Полезно придерживаться одного порядка действий.

\begin{enumerate}
\item \important{Прочитайте вопрос.} Зафиксируйте искомую величину и единицу ответа.
\item \important{Составьте «Дано».} Выпишите все числа, интервалы и скрытые связи из текста.
\item \important{Переведите слова в формулы.} Например, «составляет \(20\%\) от» означает умножение на \(0{,}2\).
\item \important{Подставьте зависимости, затем числа.} Так структура решения остаётся видимой.
\item \important{Выразите искомую величину.} Используйте переносы, пропорцию, общий знаменатель и вынесение общего множителя.
\item \important{Упростите вычисления.} Степени десяти считайте отдельно, большие числа раскладывайте на простые множители.
\item \important{Проверьте ответ.} Контролируйте знак, ограничения, единицы измерения и смысл результата.
\end{enumerate}

\important{Экзаменационная привычка:} каждый перенос и каждое действие с показателями степеней лучше записывать отдельной строкой. Такая запись предупреждает ошибки по невнимательности.

\sectiontitle{2. Перевод текста на математический язык}

\subsectiontitle{Проценты}

Если величина \(a\) составляет \(p\%\) от величины \(b\), то
\[
a=\frac{p}{100}\,b.
\]
Например, \(20\%=\frac{20}{100}=\frac15\), поэтому фраза «сила тока \(I\) составляет \(20\%\) от тока короткого замыкания \(I_{\text{кз}}\)» превращается в
\[
I=\frac15 I_{\text{кз}}.
\]

\subsectiontitle{Не менее, не более, минимум и максимум}

\begin{center}
\renewcommand{\arraystretch}{1.35}
\begin{tabularx}{\linewidth}{>{\bfseries}p{0.29\linewidth} X}
\toprule
Фраза & Математическая запись или действие \\
\midrule
Не менее \(q\) & \(x\ge q\); для минимального подходящего значения часто берут границу \(x=q\) \\
Не более \(q\) & \(x\le q\); для максимального подходящего значения часто берут границу \(x=q\) \\
\(x\) лежит от \(a\) до \(b\) & \(x\in[a;b]\), если границы допустимы \\
Наименьшее значение & Сначала выяснить направление зависимости, затем проверить нужную границу \\
Ответ в метрах & Перевести найденную длину в метры до записи ответа \\
\bottomrule
\end{tabularx}
\end{center}

\important{Читать нужно и текст, и обозначения.} Буква (L) в формуле обозначает величину, указанную в условии. Её нельзя принимать за число (1) или (21).

\sectiontitle{3. Выражение искомой величины}

\subsectiontitle{Формула сначала, числа потом}

Если дана формула
\[
P=\sigma ST^4,
\]
и требуется найти (T), удобно сначала получить общую рабочую формулу:
\[
T^4=\frac{P}{\sigma S},
\qquad
T=\sqrt[4]{\frac{P}{\sigma S}}.
\]
Только после этого следует подставлять числовые значения.

\subsectiontitle{Диагональный метод для пропорции}

Если знаменатели не равны нулю, то
\[
\frac{a}{b}=\frac{c}{d}
\quad\Longleftrightarrow\quad
ad=bc,
\qquad b\ne0,\ d\ne0.
\]
Этот приём полезен и при преобразовании формулы Доплера:
\[
f=\frac{f_0}{1-v/c}.
\]
При \(f=600\), \(f_0=590\), \(c=300\) получаем
\[
1-\frac{v}{300}=\frac{590}{600},
\]
\[
-\frac{v}{300}=\frac{590}{600}-1=-\frac{10}{600},
\]
\[
v=300\cdot\frac{10}{600}=5\ \text{м/с}.
\]

\important{Ловушка:} при переносе слагаемого знак меняется. Запись промежуточной строки защищает от случайной замены минуса плюсом.

\subsectiontitle{Уравнение и общий множитель}

Для электрической цепи пусть
\[
I=\frac{\mathcal E}{R+r},
\qquad
I_{\text{кз}}=\frac{\mathcal E}{r}.
\]
Если \(I=20\%\) от \(I_{\text{кз}}\) и \(r=1\ \text{Ом}\), то
\[
\frac{\mathcal E}{R+1}=\frac15\cdot\frac{\mathcal E}{1}.
\]
Поскольку физически \(\mathcal E\ne0\), сокращаем общий множитель \(\mathcal E\):
\[
\frac1{R+1}=\frac15,
\qquad R+1=5,
\qquad R=4\ \text{Ом}.
\]

Три регулярных вопроса:
\begin{itemize}
\item является ли полученная запись пропорцией;
\item приведено ли уравнение к удобному виду;
\item можно ли вынести или сократить общий множитель с учётом ограничений.
\end{itemize}

\sectiontitle{4. Степени и большие числа}

\subsectiontitle{Показатели степеней считаем отдельно}

Для \(a\ne0\):
\[
\frac{a^m}{a^n}=a^{m-n}.
\]
Особенно внимательно нужно обращаться с отрицательными показателями:
\[
\frac{10^{20}}{10^{-8}\cdot10^{-18}}
=10^{20-(-8)-(-18)}
=10^{46}.
\]
В вычислении полезно явно записать обе пары скобок: вычитание отрицательного числа превращается в сложение.

\subsectiontitle{Разложение на простые множители}

\[
78125=5^7,
\qquad
128=2^7,
\qquad
25=5^2,
\qquad
10=2\cdot5.
\]
Если встречается \(10^{46}\), то
\[
10^{46}=(2\cdot5)^{46}=2^{46}5^{46}.
\]
После такого разложения одинаковые основания легко объединять и сокращать.

\subsectiontitle{Извлечение корня}

Если \(T^4=A\) и по смыслу \(T>0\), то
\[
T=A^{1/4}.
\]
Даже если показатели после умножения на \(\frac14\) получаются дробными, запись остаётся допустимой. Необычный вид ответа служит поводом ещё раз проверить исходные данные и вычисления.

\sectiontitle{5. Анализ зависимости величин}

Рассмотрим формулу тонкой линзы:
\[
\frac1{d_1}+\frac1{d_2}=\frac1f,
\]
где \(f=30\ \text{см}\), \(d_1\in[30;50]\), \(d_2\in[180;210]\). Требуется найти наименьшее возможное \(d_1\).

Выразим зависимость:
\[
\frac1{d_1}=\frac1f-\frac1{d_2}.
\]

\begin{enumerate}
\item Чтобы \(d_1\) было наименьшим, дробь \(\frac1{d_1}\) должна быть наибольшей.
\item Величина \(\frac1f\) постоянна.
\item Чтобы разность была наибольшей, вычитаемая дробь \(\frac1{d_2}\) должна быть наименьшей.
\item Дробь \(\frac1{d_2}\) наименьшая при наибольшем \(d_2\), то есть при \(d_2=210\).
\end{enumerate}

Подставляем:
\[
\frac1{d_1}=\frac1{30}-\frac1{210}
=\frac7{210}-\frac1{210}
=\frac6{210}=\frac1{35}.
\]
Следовательно,
\[
d_1=35\ \text{см}.
\]
Число (35) принадлежит заданному интервалу ([30;50]), поэтому ограничение выполнено.

\important{Общий принцип:} для положительных (x) величины (x) и (1/x) меняются в противоположных направлениях.

\sectiontitle{6. Размерности и смысл ответа}

Формула обычно возвращает величину в тех единицах, в которых записаны исходные данные. Если найдено \(h=0{,}00125\ \text{км}\), а ответ нужен в метрах, то
\[
0{,}00125\ \text{км}=0{,}00125\cdot1000=1{,}25\ \text{м}.
\]

Перед финальной записью проверьте:
\begin{itemize}
\item положительна ли физическая величина;
\item принадлежит ли найденное значение заданному интервалу;
\item не обращается ли знаменатель в ноль;
\item выполнена ли исходная формула после подстановки;
\item совпадают ли требуемые и полученные единицы;
\item реалистичен ли порядок результата.
\end{itemize}

\sectiontitle{7. Типичные ошибки}

\begin{enumerate}
\item Пропустить скрытое данное: процент, интервал, слово «минимальный» или обозначение (L).
\item Сразу подставлять числа и потерять структуру формулы.
\item Поменять знак неверно при переносе слагаемого.
\item Считать (20-(-8)-(-18)) устно и ошибиться в знаках.
\item Не разложить \(10^n\) как \(2^n5^n\), когда это нужно для сокращения.
\item Сократить множитель, не проверив, может ли он равняться нулю.
\item Выбрать неправильную границу интервала при поиске минимума.
\item Записать ответ в километрах, когда требуются метры.
\end{enumerate}

\sectiontitle{8. Итоговый чек-лист}

Перед тем как завершить задачу, отметьте каждый пункт:
\begin{itemize}[label=$\square$]
\item Я выписала искомую величину и единицу ответа.
\item Я перенесла из текста все числа, интервалы и смысловые ограничения.
\item Я перевела проценты и словесные связи в формулы.
\item Я сначала подставила формулы, затем числовые значения.
\item Я записала действия со знаками и степенями отдельными строками.
\item Я проверила возможность сокращения и ограничения знаменателей.
\item Я проанализировала направление зависимости при поиске минимума или максимума.
\item Я проверила единицы измерения и смысл ответа.
\end{itemize}

\newpage
\sectiontitle{Тренировочный блок}

\textbf{Средний уровень}
\begin{enumerate}[label=\textbf{\arabic*.}]
\item По формуле \(F=ma\) найдите \(F\), если \(m=2{,}4\ \text{кг}\), \(a=3{,}5\ \text{м/с}^2\).
\item Величина \(a\) составляет \(20\%\) от \(450\). Найдите \(a\).
\item Мощность задаётся формулой \(P=UI\). Найдите \(I\), если \(P=1320\ \text{Вт}\), \(U=220\ \text{В}\).
\end{enumerate}

\textbf{Уровень выше среднего}
\begin{enumerate}[label=\textbf{\arabic*.},start=4]
\item Для тонкой линзы \(\frac1{d_1}+\frac1{d_2}=\frac1f\). Известно: \(f=24\ \text{см}\), \(d_2\in[120;168]\ \text{см}\). Найдите наименьшее \(d_1\).
\item Частота гудка задаётся формулой \(f=\frac{f_0}{1-v/c}\). Найдите минимальную скорость \(v\), если \(f_0=590\ \text{Гц}\), различимая частота должна быть не менее \(600\ \text{Гц}\), \(c=300\ \text{м/с}\).
\item Даны \(I=\frac{\mathcal E}{R+r}\), \(I_{\text{кз}}=\frac{\mathcal E}{r}\). Найдите \(R\), если \(I=25\%\) от \(I_{\text{кз}}\), \(r=2\ \text{Ом}\), \(\mathcal E\ne0\).
\item По формуле \(P=\sigma ST^4\) найдите положительную температуру \(T\), если \(P=567\), \(\sigma=5{,}67\cdot10^{-8}\), \(S=10^{-2}\).
\item Упростите \(\dfrac{10^{17}}{10^{-5}\cdot10^8}\).
\item Найдите положительное \(x\), если \(x^4=2^{20}\cdot5^{12}\).
\end{enumerate}

\textbf{Сложный уровень}
\begin{enumerate}[label=\textbf{\arabic*.},start=10]
\item Для положительных величин выполняется
\[
Q=\frac{kx^4}{y},
\]
где \(Q=2{,}5\cdot10^6\), \(k=4\cdot10^{-3}\), \(y=10^2\). Найдите \(x\) и проверьте результат подстановкой.
\end{enumerate}

\newpage
\sectiontitle{Ответы к тренировочному блоку}

\begin{center}
\renewcommand{\arraystretch}{1.45}
\begin{tabularx}{\linewidth}{c X}
\toprule
№ & Ответ \\
\midrule
1 & \(8{,}4\ \text{Н}\) \\
2 & (90) \\
3 & \(6\ \text{А}\) \\
4 & \(28\ \text{см}\) \\
5 & \(5\ \text{м/с}\) \\
6 & \(6\ \text{Ом}\) \\
7 & (1000) \\
8 & \(10^{14}\) \\
9 & \(x=2^5\cdot5^3=4000\) \\
10 & \(x=500\) \\
\bottomrule
\end{tabularx}
\end{center}

\vfill
\begin{center}
\textcolor{softgray}{Код самопроверки: ФОРМУЛА--ЗНАК--ЕДИНИЦА}
\end{center}

\end{document}
